% !Mode:: "TeX:UTF-8"

\section{项目设计目标}

随着物联网技术和可穿戴设备的普及，传统手表的功能已无法满足现代消费者对健康监测、运动记录和数据交互的需求。智能手表作为一种集成度高、功能丰富的可穿戴设备，代表了嵌入式系统在消费电子领域的重要应用方向。基于STM32微控制器的智能手表系统，能够通过先进的传感器融合技术和低功耗通信技术实现更人性化、智能化的用户体验。

STM32F103C8T6微控制器凭借其72MHz Cortex-M3内核、丰富的外设接口（I2C、USART、定时器、DMA等）以及低功耗特性，为可穿戴设备开发提供了理想的解决方案。本项目结合嵌入式系统课程所学知识，选择STM32F103C8T6作为核心控制器，设计并实现一款集成OLED显示、六轴运动传感、蓝牙通信与计步功能的智能手表系统，旨在将理论知识与实际应用相结合，探索嵌入式系统在可穿戴设备领域的应用潜力。

本项目旨在设计并实现一个功能完备、运行稳定的智能手表系统，主要目标包括：
\begin{itemize}
    \item \textbf{实现OLED图形显示系统：} 基于I2C总线驱动SSD1306显示控制器，开发支持16$\times$24大字时钟、8$\times$16标准字符、像素级绘图以及多页面循环切换的UI框架。
    \item \textbf{开发六轴运动传感与姿态解算：} 通过MPU6050采集加速度和陀螺仪数据，实现加速度量程$\pm$2g（16384 LSB/g）与陀螺仪量程$\pm$250°/s（131 LSB/°/s）的物理量转换，并计算Pitch和Roll姿态角。
    \item \textbf{实现HC-05蓝牙通信协议：} 设计基于USART2 + DMA的环形缓冲接收机制，定义自定义帧协议（STX=0xAA, ETX=0x55, XOR校验），支持传感器数据上报、时间同步、计步数据传输三种命令类型。
    \item \textbf{设计计步器算法：} 基于加速度向量幅值与低通滤波，配合动态阈值和200ms最小步间间隔防抖，实现步数统计、距离估算和卡路里消耗计算。
    \item \textbf{开发Android上位机应用：} 实现蓝牙设备扫描、连接管理、传感器数据实时展示、时间同步下发、通信日志记录等功能，提供美观的Material Design 3用户界面。
    \item \textbf{完成系统集成与原型验证：} 在STM32F103C8T6最小系统板及外围模块搭建的硬件平台上进行功能测试与性能评估，确保系统稳定运行并通过I2C总线恢复机制增强鲁棒性。
\end{itemize}